\documentclass[11pt,letterpaper]{article}

\usepackage[margin=1in]{geometry}
\usepackage{termcal}
\usepackage{enumitem}
\usepackage[colorlinks=true, allcolors=blue]{hyperref}
\usepackage{color}
\usepackage{multirow}
\usepackage{multicol}

\newcommand{\todo}[1]{\textcolor{red}{TODO: #1}}

\title{16.07: Dynamics}
\author{Fall 2026}
\date{}

\begin{document}

\maketitle

%\noindent 
%\textbf{Course website:} \href{https://optimalcontrol.ri.cmu.edu/}{optimalcontrol.ri.cmu.edu}

\section*{Course Description}

This course is about modeling and simulating mechanical systems, starting with a review of basic Newtonian mechanics of point-mass particles, then covering rigid-body and multi-body systems. Applications to aircraft, spacecraft, and robotic systems will be emphasized.

\medskip
\noindent
\textbf{Prerequisites:} Officially -- (16.001 or 16.002) and (16.003 or 16.004). Practically -- basic mechanics at the level of 8.01, calculus and vectors/matrices at the level of 18.02, and prior experience with a high-level programming language like Python or MATLAB (we'll use MATLAB).

\section*{Instructors}

\begin{center}
\begin{tabular}{l l l}
	Lecturer: Zac Manchester & \textbf{Email:} \href{mailto:zacm@mit.edu}{zacm@mit.edu} & \textbf{Office:} 31-243 \\
	TA: John Zhang & \textbf{Email:} \href{mailto:jzhang3@mit.edu}{jzhang3@mit.edu}
\end{tabular}
\end{center}

\section*{Logistics}

\begin{itemize}
	\item Lectures will be held \textbf{Tuesdays and Thursdays 11:00--12:30 in 35-225}. Lectures will also be recorded for later viewing.
	\item Recitations and labs will be held \textbf{Fridays 3:00--4:00 in \todo{TBD}}.
	\item Office hours will be \todo{TBD}.
	\item Homework assignments will be due \textbf{Fridays by 5:00 PM}.
	\item There will be occasional short (5 minute) quizzes given in class.
	\item There will be an in-class mid-term exam and a cumulative final exam.
	\item Piazza will be used for general discussion and Q\&A outside of class and office hours.
\end{itemize}

\newpage 
\section*{Learning Objectives}
By the end of this course, students should be able to do the following:
\begin{enumerate}
	\item Select appropriate coordinates to describe the motion of particle and multi-body system
	\item Formulate dynamics in non-inertial reference frames
	\item Identify and exploit conserved quantities like energy and momentum
	\item Derive equations of motion for multi-body systems using the Euler-Lagrange equation
	\item Model the orbital dynamics of a spacecraft as central-force motion of a particle
	\item Use Euler's equation to model and analyze the attitude dynamics of a rigid body and apply this to aircraft and spacecraft
\end{enumerate}

\section*{Learning Resources}

There is no required textbook for this course. Video recordings of lectures and lecture notes will be posted online. The following books are good references for the topics discussed in the course that you may want to refer to:

\begin{enumerate}
\item J. L. Meriam, L. G. Kraige, and J. N. Bolton, Engineering Mechanics: Dynamics, 9th ed. Hoboken, NJ, USA: Wiley, 2020.
	\item L. D. Landau and E. M. Lifshitz, Mechanics, 3rd ed. Oxford, U.K.: Butterworth-Heinemann, 1976.
	\item H. Goldstein, C. Poole, and J. Safko, Classical Mechanics, 3rd ed. San Francisco, CA, USA: Addison Wesley, 2002.
	\item K. M. Lynch and F. C. Park, Modern Robotics: Mechanics, Planning, and Control. Cambridge, U.K.: Cambridge University Press, 2017.
\end{enumerate}

\section*{Assignments}

\begin{itemize}
	\item Students will be asked to complete homework assignments roughly weekly that include both an analytical question and a numerical simulation question to be completed in MATLAB.
	\item There will be three labs throughout the semester during the Friday recitation time. On lab weeks, homework assignments will include analysis of experiment data.
\end{itemize}

\section*{Quizzes and Exams}

\begin{itemize}
	\item Short quizzes will be given approximately weekly in class. These are intended to take no more than 5 minutes to complete.
	\item There will be an in-class mid-term exam and a cumulative final exam.
\end{itemize}

\section*{Grading}

Grading will be based on:
\begin{itemize}
	\item 25\% Homework and Labs
	\item 15\% Quizes
	\item 25\% Mid-Term Exam
	\item 30\% Final Exam
	\item 5\% Participation
\end{itemize}
Students can earn a full participation grade through any combination of office hours, Piazza discussions, recitation attendance, and by offering constructive feedback about the course to the instructors.


\section*{Course Policies}

\textbf{Late Assignments:} Students are allowed a budget of six late days for turning in homework with no penalty throughout the semester. They may be used together on one assignment, or separately across multiple assignments. Beyond these six days, no other late homework will be accepted.

\noindent \textbf{Collaboration:} Students are encouraged to collaborate on labs and homework assignments, but each student must turn in their own assignment by the deadline.

\noindent \textbf{AI Use:} Students are permitted to use AI tools like Claude and Gemini for assistance on homework assignments. This should be acknowledged explicitly with a brief statement about how AI was used.

\section*{Tentative Schedule}

\begin{tabular}{c|c|c|c}
	Week & Dates & Topics & Assignments \\
	\hline
	\multirow{2}{*}{1} & Sep 10 & 
        Course Overview, History, and Review & Survey \\
	 &  &  &  \\
	\hline
	\multirow{2}{*}{2} & Sep 15 &
        Coordinates and Reference Frames &  \\
	 & Sep 17 & Momentum, Energy, and Conservation & HW1 Out \\
	\hline
	\multirow{2}{*}{3} & Sep 22 &
        Euler-Lagrange Equations &  \\
	 & Sep 24 & Potential Functions, Non-Conservative Forces &  \\
	\hline
	\multirow{2}{*}{4}  & Sep 29 &
        Simulation and Runge-Kutta \textcolor{red}{(Zac Out)} & HW1 Due \\
	 & Oct 1 & Constraints and Lagrange Multipliers & HW2 Out \\
	\hline
	\multirow{2}{*}{5}  & Oct 6 & Orbit Mechanics 1 &  \\
	 & Oct 8 & Orbit Mechanics 2 &  \\
	\hline
	\multirow{2}{*}{6}  & Oct 13 & \textcolor{red}{No Class --- Monday Schedule} & HW2 Due \\
	 & Oct 15 & Rigid Bodies and Attitude & HW3 Out \\
	\hline
	\multirow{2}{*}{7}  & Oct 20 & Rotation Matrices and Quaternions 
         &  \\
	 & Oct 22 & Euler's Equation & \\
	\hline
	\multirow{2}{*}{8}  & Oct 27 & 
        Newton-Euler Dynamics for Aircraft & \\
	 & Oct 29 & \textbf{Mid-Term Exam} \textcolor{red}{(Zac Out)} &   \\
	\hline
	\multirow{2}{*}{9}  & Nov 3 & Linearization and Stability \textcolor{red}{(Zac Out)}         &  \\
	 & Nov 5 & Linear Dynamical Systems and Modes \textcolor{red}{(Zac Out)} &  \\
	\hline
	\multirow{2}{*}{10}  & Nov 10 & Vector Fields and Phase Portraits
         & HW3 Due \\
	 & Nov 12 & Coulomb Friction and the Maximum Dissipation Principle & HW4 Out  \\
	 \hline
	\multirow{2}{*}{11}  & Nov 17 & 3D Multi-Body Dynamics Simulation
         &  \\
	 & Nov 19 & Constraint Stabilization &   \\
	 \hline
	\multirow{2}{*}{12}  & Nov 24 & System Identification and Least-Squares
         &  HW4 Due \\
	 & Nov 26 & \textcolor{red}{No Class --- Thanksgiving} & HW5 Out  \\
	 \hline
	\multirow{2}{*}{13}  & Dec 1 & TBD: Schedule Margin
         &  \\
	 & Dec 3 &  &   \\
	 \hline
	\multirow{2}{*}{14}  & Dec 8 & TBD: Schedule Margin
         & HW5 Due \\
	 & Dec 10 &  &   \\

\end{tabular}


\end{document}
